% Copyright 2026 Sewoong Lee
% 
% This work may be distributed and/or modified under the conditions of the
% LaTeX Project Public License, either version 1.3c of this license or (at your
% option) any later version. The latest version of this license is at:
% https://www.latex-project.org/lppl.txt
% 
% This work has the LPPL maintenance status `maintained'.
% The Current Maintainer is Sewoong Lee <samuel27@illinois.edu>.
% 
% The work consists of README.md, LICENSE, beamercolorthemechalkboard.sty,
% beamercolortheme-chalkboard-doc.tex, and chalkboard-example.tex.
% The derived files are beamercolortheme-chalkboard-doc.pdf and chalkboard-example.pdf.

\documentclass[aspectratio=169,11pt]{beamer}
\usetheme{Madrid}
\usecolortheme{chalkboard}
\title[Chalkboard]{Chalkboard}
\subtitle{A dark green color theme for Beamer}
\author{Sewoong Lee}
\date{Version 0.1.0 -- October 15, 2026}
\begin{document}
\begin{frame}
\titlepage
\end{frame}
\section{Examples}
\begin{frame}{Text, lists, and mathematics}
Warm off-white text on a dark green board.
\begin{itemize}
\item Structure uses chalk yellow.
\item \alert{Alerts use chalk pink.}
\item {\usebeamercolor[fg]{example text}Examples use chalk blue.}
\end{itemize}
\begin{enumerate}
\item Start with a clear statement.
\item Support it with a simple equation:
\end{enumerate}
\[ e^{i\pi}+1=0 \qquad \sum_{k=1}^{n}k=\frac{n(n+1)}{2}. \]
\end{frame}
\begin{frame}{Three kinds of blocks}
\begin{block}{A standard block}
Use this for definitions and central ideas.
\end{block}
\begin{alertblock}{An alert block}
Use an alert when a detail deserves special attention.
\end{alertblock}
\begin{exampleblock}{An example block}
For example, $a^2+b^2=c^2$ for a right triangle.
\end{exampleblock}
\end{frame}
\begin{frame}{The chalk palette}
\centering
\renewcommand{\arraystretch}{1.4}
\begin{tabular}{lll}
\textbf{Color} & \textbf{HTML} & \textbf{Sample}\\
chalkwhite & E6E4D8 & \textcolor{chalkwhite}{Chalk text}\\
chalkyellow & EBD58A & \textcolor{chalkyellow}{Structure}\\
chalkblue & 9FC5E0 & \textcolor{chalkblue}{Examples}\\
chalkpink & F0A6A0 & \textcolor{chalkpink}{Alerts}\\
chalkboard & 2A3B33 & Background\\
chalkboarddark & 1F2B25 & Titles and palettes\\
chalkboardlight & 34483E & Block bodies
\end{tabular}
\end{frame}
\end{document}
